% \vspace{-10pt}
\section{Memory-Augmented Manipulation Policies}
\vspace{4pt}
Building on \benchname, we construct a family of memory-augmented vision-language-action (VLA) models based on the $\pi_{0.5}$ backbone, collectively termed the \textbf{\modelname\xspace suite}. 
We systematically compare different memory representations and integration mechanisms under controlled settings, as illustrated in Figure~\ref{fig:memory_design}.
More detailed model formulations are provided in Appendix~\ref{sec:appx_model_design}.
% $\pi_{0.5}$ is based on PaLI-Gemma \cite{beyer2024paligemma}, using SigLIP2 \cite{tschannen2025siglip} for vision encoding and Gemma \cite{team2024gemma} as the language backbone.
% $\pi_{0.5}$ encodes multi-view images into image tokens and concatenates them with language tokens for task instructions and proprioceptive states to an vision-language-model (VLM) expert, then predicts actions with a flow-matching-based action expert. 
%(need to find a way for naming)
\subsection{Memory Representations}

We study three representations: \textit{symbolic}, \textit{perceptual}, and \textit{recurrent}. 
Symbolic memory uses discrete language tokens, while the latter two provide differentiable neural features.

\noindent{\bf Symbolic Memory} is represented as interpretable language subgoals, as correct subgoals imply effective history summarization. At each step, we leverage an auxiliary vision-language model (VLM) that conditions on the current image and prior subgoals to generate the next subgoal. We consider \emph{simple} instructions and \emph{grounded} variants that additionally include image coordinates (e.g., \textit{pick up the green cube} vs. \textit{pick up the green cube at [63, 152]}), with grounding being particularly beneficial for spatial reasoning. Notably, symbolic memory approaches typically require additional subgoal annotations to train the VLM.

\noindent{\bf Perceptual Memory} is represented as a sequence of visual tokens selected from past images and  extracted by the $\pi_{0.5}$ vision encoder. To select informative tokens, we adopt two strategies: (1) token dropping \cite{yao2025timechat}, which removes temporally redundant image patches based on RGB differences; (2) uniform frame sampling \cite{chen2024videollm}, which evenly downsamples the sequence and concatenates tokens from the sampled frames. In practice, we found that visual tokens alone suffice to encode history without proprioceptive states.

\noindent{\bf Recurrent Memory} compresses the visual token sequence into fixed-size latent states through recurrence. We adopt two models: (1) Test-Time Training (TTT) \cite{sun2024learning, zhang2025test}, which updates fast weights online via a self-supervised loss and applies them to generate output features; and (2) Recurrent Memory Transformers (RMT) \cite{bulatov2022recurrent, bulatov2023scaling}, which process the input sequence into segments and recurrently update a set of learnable memory slots for each segment using a transformer \cite{vaswani2017attention}.

% \noindent{\bf Symbolic Memory (SM)} represents memory through interpretable language subgoals generated by an auxiliary vision-language model (VLM). 
% At each step, the VLM predicts the next subgoal conditioned on the current image and prior subgoals, which together form an explicit memory of past interactions. We study two variants: \emph{simple} subgoals that provide brief instructions, and \emph{grounded} subgoals that additionally include image coordinates (e.g., \textit{pick up the green cube"} vs. \textit{pick up the green cube at [63, 152]"}). As shown in our experiments, grounding is particularly important for tasks requiring precise spatial reasoning.

% \noindent{\bf Perceptual Memory (PM)} stores visual tokens from past images extracted by SigLIP2 \cite{tschannen2025siglip}, the vision encoder of $\pi_{0.5}$. Retaining all historical tokens would excessively increase the context length and slow inference, so we select informative tokens using two common strategies from streaming video understanding: (1) token dropping \cite{yao2025timechat}, which removes temporally redundant tokens based on RGB differences, and (2) uniform frame sampling \cite{chen2024videollm}, which downsamples long sequences and concatenates tokens from sampled frames as memory. Although proprioceptive states can be incorporated for state-aware features, we find that image tokens alone suffice to encode history.

% \noindent{\bf Recurrent Memory (RM)} compresses historical visual tokens into fixed-size latent states through recurrence. We study two recurrent architectures: (1) Test-Time Training (TTT) \cite{sun2024learning, zhang2025test}, which sequentially updates tokens via online fast-weight optimization and uses the resulting fast-weight projections as memory; and (2) Recurrent Memory Transformers (RMT) \cite{bulatov2022recurrent, bulatov2023scaling, wang2025videollamb}, which handle long contexts by chunking inputs and recurrently updating a fixed set of learnable tokens through attention, enabling stable long-horizon storage and retrieval.


\subsection{Memory Integration Mechanisms}
Symbolic memory can be naturally integrated into $\pi_{0.5}$  by concatenating subgoals with task instructions. In contrast, perceptual and recurrent memory produce neural \emph{memory tokens} that require more architectural modifications. We therefore study three integration mechanisms.


\noindent\textbf{Memory-as-Context.}
Memory tokens are concatenated with the original inputs and jointly processed by the VLM expert, directly influencing VLM features.

\noindent\textbf{Memory-as-Modulator.}
We adopt adaptive LayerNorm (AdaLN) \cite{peebles2023scalable} to condition the action expert on external memory. Before entering the feed-forward block in each layer, the action features cross-attend to memory tokens via multi-head attention to extract memory-aware representations. These representations are then projected into scale and shift parameters that modulate the normalized action features through AdaLN.
% For simplicity, the self-attention layers of the action expert remain unchanged.
% The attention layers of the action expert remain unchanged for simplicity.

\noindent\textbf{Memory-as-Expert.}
We introduce an additional lightweight memory expert that processes memory tokens through a dedicated pathway. The three experts interact via a specified blockwise causal attention scheme: the action expert attends to both the VLM and memory experts, while the VLM and memory experts do not attend to each other. This separation limits interference, preserves the original VLM behavior, and allocates specialized capacity for memory processing.


% \noindent\textbf{Memory-as-Context.}
% We concatenate memory tokens with the original input tokens and feed them directly into the VLM expert in $\pi_{0.5}$ for fine-tuning. Because the VLM uses bidirectional attention, memory tokens influence its internal representations at all layers.

% \noindent\textbf{Memory-as-Modulation.}
% Inspired by Diffusion Transformers \cite{peebles2023scalable}, we apply adaptive LayerNorm Zero (adaLN-Zero) to condition the action expert through learned scale and shift parameters. At each feed-forward layer, action features cross-attend to memory tokens via a lightweight transformer to obtain memory-aware features, which are then mapped to modulation parameters that adapt the action activations. Attention layers in the action expert remain unchanged for simplicity.

% \noindent\textbf{Memory-as-Expert.}
% We add a lightweight \emph{memory expert} that processes memory tokens in a separate pathway. All experts communicate through block-wise causal attention: tokens attend bidirectionally within each block but not across future blocks. The action expert attends to both VLM expert and memory outputs, while the other experts perform self-attention only. 
% This design preserves the original VLM behavior and limits cross-expert interference.

